\chapter{Réalisation logicielle}
\label{chap:realisation}

\section{Architecture d'ensemble}
La figure~\ref{fig:systeme} présente la chaîne complète, des sources de données au navigateur de
l'utilisateur. Le calcul se fait là où sont le GPU et les poids (Google Colab), tandis que
l'adresse publique reste fixe grâce à un Worker Cloudflare et à un tunnel nommé.

\begin{figure}[H]
\centering
\begin{tikzpicture}[
  b/.style={draw, rounded corners=3pt, align=center, font=\footnotesize, minimum height=1.1cm,
            text width=3.3cm, fill=bleuespirt!8},
  s/.style={b, fill=gray!10},
  g/.style={b, fill=orange!12},
  fl/.style={-{Stealth[length=2mm]}, thick}]
  \node[b] (nav) at (0, 0) {Navigateur\\application ATOM Wind};
  \node[b] (worker) at (4.4, 0) {Worker Cloudflare\\atom.zelqane.com\\(clé injectée, cache)};
  \node[b] (tunnel) at (8.8, 0) {Tunnel Cloudflare\\nommé atom-origin};
  \node[s] (drive) at (0, -2.3) {Google Drive\\poids, cache, historique};
  \node[g] (modele) at (4.4, -2.3) {Modèle PC-FNO3D\\(GPU)};
  \node[g] (service) at (8.8, -2.3) {Service FastAPI\\et serveur MCP\\(Google Colab)};
  \node[s] (client) at (4.4, -4.6) {Clients de l'API\\et agents MCP};
  \node[s] (src) at (8.8, -4.6) {GFS (Open-Meteo),\\METAR, SYNOP};
  \draw[fl] (nav) -- (worker);
  \draw[fl] (worker) -- (tunnel);
  \draw[fl] (tunnel) -- (service);
  \draw[fl] (service) -- (modele);
  \draw[fl] (drive) -- (modele);
  \draw[fl] (src) -- (service);
  \draw[fl] (client) -- (service);
\end{tikzpicture}
\caption{Architecture de déploiement.}
\label{fig:systeme}
\end{figure}

\section{Le paquet \texttt{pcfno3d}}
Le notebook de recherche a été porté bloc par bloc en un paquet Python installable (41 modules,
environ 16\,800 lignes), qui est le seul chemin de code utilisé en production. Chaque bloc du
notebook devient un module ; la configuration et la géométrie sont passées explicitement au lieu
de variables globales (tableau~\ref{tab:modules}). Le portage a été validé en retrouvant
exactement la table de résultats du notebook.

\begin{table}[H]
\centering
\caption{Principaux modules du paquet.}
\label{tab:modules}
\begin{tabularx}{\linewidth}{@{}lX@{}}
\toprule
Module & Rôle\\
\midrule
\code{config}, \code{geometry} & Configuration typée ; grille, relief, masques, opérateurs discrets (gradient, résidu de masse)\\
\code{contract}, \code{cache} & Contrat de données ANEO ; tenseurs par instant et index\\
\code{dataset}, \code{features} & Entrée du modèle : canaux, vent géostrophique ajusté, entrée de production\\
\code{model}, \code{losses}, \code{train} & PC-FNO3D, perte physique, entraînement (EMA, OneCycle, reprise)\\
\code{evaluate}, \code{dso} & Métriques cumulées, verdicts DSO1 à DSO6\\
\code{inference}, \code{reliability} & Modèle de production, ensemble d'incertitude différé\\
\code{forcage}, \code{temps\_reel}, \code{gfs} & Forçage ERA5, temps réel GFS, archive GFS (OS7)\\
\code{alerts}, \code{windfarm}, \code{publish} & Alertes, post-traitement éolien, publication et exports NetCDF/CSV\\
\code{os6}, \code{atlas}, \code{baselines} & Plan du réseau, atlas éolien, méthodes de référence\\
\bottomrule
\end{tabularx}
\end{table}

Une interface en ligne de commande (\code{python -m pcfno3d <commande>}) expose toutes les étapes :
construction du cache, contrôle, entraînement, évaluation, reconstruction, alertes, publication,
service. Chaque commande accepte une configuration et des surcharges \code{--set section.clé=valeur}.

\section{Service web et sécurité}

\subsection{API REST}
Le service est une application FastAPI. Ses routes donnent la grille, les journées calculées, le
vent en un point et un profil vertical, la reconstruction à partir de mesures envoyées (fichier CSV
validé par le contrat), les exports NetCDF et CSV, la qualité du modèle servi, les alertes,
l'historique et les espaces métier (agriculture, production éolienne, maintenance, dimensionnement
de mât, espaces holding, chantier et environnement).

\subsection{Sécurité}
\begin{itemize}
  \item \textbf{fermé par défaut} : sans clé, toutes les routes répondent 401, sauf l'état de santé ;
  \item \textbf{clés d'API} de la forme \code{pcf\_<32 caractères>}, stockées uniquement par leur
        empreinte SHA-256, avec deux portées (lecture, administration), un quota par clé (réponse
        429) et un journal d'accès qui ne contient jamais de secret ; une révocation prend effet à la
        requête suivante ;
  \item \textbf{comptes utilisateurs} : mots de passe salés et dérivés par scrypt, blocage après
        cinq échecs, trois rôles (agriculteur, exploitant éolien validé par l'administrateur,
        administrateur) vérifiés côté serveur à chaque requête ;
  \item \textbf{secrets} : la clé du Worker n'est connue du service que par son empreinte, et les
        identifiants du tunnel ne transitent que par une variable d'environnement, jamais par le
        dépôt.
\end{itemize}

\subsection{Serveur MCP et assistant}
Le même modèle est exposé par un serveur MCP (\emph{Model Context Protocol}) de 22 outils :
consultation du vent et des profils, prévision en un point, décision agricole, fenêtres de
maintenance, valeur et influence des stations, simulation de panne, historique\dots{} Un agent
conversationnel peut ainsi lancer une simulation et en lire le résultat sans intervention (critère
MCP1). L'application intègre aussi un assistant fondé sur Claude, qui répond aux questions des
utilisateurs en appelant ces mêmes outils.

\section{Temps réel et prévision}
La réanalyse ERA5 arrive avec cinq jours de retard ; le temps réel utilise donc la prévision GFS
(via l'API Open-Meteo), calculée sur les mêmes mailles et avec les mêmes formules que le forçage
d'entraînement, et les mesures du jour : METAR de Tunis-Carthage et SYNOP de Zaghouan et Nabeul.
Le service calcule la journée en cours (heures passées avec les stations, heures futures en
prévision), le lendemain et le surlendemain, les garde en cache trente minutes et les archive sur
disque, ce qui permet de comparer plus tard une prévision au temps réel du même jour. Un ensemble
GEFS de 31 membres donne la probabilité de pouvoir traiter.

\section{L'application ATOM Wind}
L'application, construite à partir de la maquette fonctionnelle conçue pendant le projet, est une
page web sans étape de compilation, branchée sur l'API.

\begin{itemize}
  \item \textbf{Mon vent} : vent à 10\,m et à 2\,m, rafales, direction, pour un lieu choisi sur la
        carte, heure par heure ;
  \item \textbf{espace agriculteur} (application principale) : fenêtres de traitement et
        d'aspersion de J0 à J+2 avec la règle appliquée affichée, besoin en eau par la méthode
        FAO-56 \citep{allen1998} avec la température, l'humidité et le rayonnement de GFS ;
  \item \textbf{démonstration éolienne} : vent au moyeu et sur le rotor, production d'éoliennes de
        référence avec effets de sillage \citep{jensen1983}, rose des vents, carte du gisement et
        atlas annuel ;
  \item \textbf{vue 3D} : relief exagéré, lignes de courant, coupe horizontale, couleur selon la
        vitesse ou l'incertitude, colonnes de vent par nœud, relecture de la journée, fiche d'un
        point avec son profil et son intervalle à 80\,\% ;
  \item \textbf{tableau de bord} : indicateurs du service (stations actives, temps de
        reconstruction, résidu de masse de l'instant) et qualité du modèle servi ;
  \item \textbf{alertes} : station muette, résidu de masse hors seuil, vent fort, avec
        notifications par courriel, SMS ou WhatsApp ;
  \item \textbf{interface trilingue} français, arabe (de droite à gauche) et anglais ;
  \item \textbf{continuité de service} : quand le service est arrêté, le Worker sert le dernier
        état réel connu, signalé comme tel.
\end{itemize}

\acompleter{Insérer ici 3 à 5 captures d'écran de l'application (Mon vent, espace agriculteur,
vue 3D, tableau de bord) : déposer les images dans figures/ et utiliser \textbackslash includegraphics.}

\section{Qualité logicielle}
\begin{itemize}
  \item \textbf{tests} : 435 fonctions de test dans 42 fichiers, exécutées sur CPU en quelques
        minutes sur un mini-jeu de données synthétisé au format exact du générateur ; elles
        couvrent la physique (résidu de masse, continuité), l'API (sécurité, quotas), le contrat de
        données, la publication (atomicité face à un serveur simulé avec coupures), l'interface
        (dictionnaires de traduction) et le service Colab ;
  \item \textbf{reproductibilité} : même entrée et même version donnent la même sortie au bit près
        sur CPU ; sur GPU, l'écart reste de quelques $10^{-3}$\ms (calcul en bfloat16) ;
  \item \textbf{traçabilité} : chaque export porte l'empreinte des poids, des statistiques de
        normalisation et l'horodatage ; le cache porte un manifeste vérifiable ;
  \item \textbf{documentation} : fiches du modèle et des données, document d'intégration pour les
        équipes qui consomment la sortie, diagrammes UML.
\end{itemize}
